\documentclass[article,12pt,oneside,a4paper,english,brazil,sumario=tradicional]{abntex2}
\usepackage[T1]{fontenc}
\usepackage[utf8]{inputenc}
\usepackage{graphicx}
\usepackage[alf]{abntex2cite}

\titulo{Revisão automatizada de manuscritos científicos em português}
\autor{Fulana de Tal}
\data{2026}

\begin{document}
\maketitle

\begin{resumoumacoluna}
Este trabalho apresenta uma ferramenta de linha de comando que audita artigos
científicos escritos em LaTeX. A ferramenta verifica estrutura, referências,
normas da ABNT e qualidade da escrita, e aponta indícios que merecem revisão
humana. Os resultados em um artigo de exemplo mostram que os problemas
plantados são encontrados com a linha correta, o que facilita a correção
antes da submissão a um periódico.

\vspace{\onelineskip}
\noindent
\textbf{Palavras-chave}: revisão automatizada. processamento de linguagem natural. escrita acadêmica.
\end{resumoumacoluna}

\renewcommand{\resumoname}{Abstract}
\begin{otherlanguage*}{english}
\begin{resumoumacoluna}
This paper presents a command-line tool that audits scientific papers written
in LaTeX. The tool checks structure, references, ABNT standards and writing
quality, and points out signals that deserve human review. Results on a sample
paper show that the planted problems are found at the correct line, which helps
authors fix them before submitting to a journal.

\vspace{\onelineskip}
\noindent
\textbf{Keywords}: automated review. natural language processing. academic writing.
\end{resumoumacoluna}
\end{otherlanguage*}

\textual
\section{Introdução}\label{sec:introducao}

Ferramentas de PLN têm sido aplicadas à revisão de textos acadêmicos há pelo
menos uma década \cite{lecun2015}. Neste trabalho, o processamento de linguagem
natural (PLN) é usado para apoiar autores que escrevem em português.

A revisão manual de manuscritos é custosa devido à fatores como o volume
crescente de submissões. À partir de 2020, diversos periódicos passaram a
exigir verificações automáticas de similaridade \cite{devlin2019}.

A idéia central deste trabalho é avaliar o texto a nível de parágrafo e de
seção, apontando problemas que um revisor humano costuma encontrar.

O objetivo deste trabalho é descrever uma ferramenta de auditoria que lê os
arquivos-fonte de um artigo, identifica problemas de estrutura, de referências
e de escrita, e apresenta cada achado com o arquivo, a linha e uma sugestão de
correção, sem jamais alterar os arquivos originais do autor.

Trabalhos anteriores \cite{silva2099} propuseram verificadores para a língua
inglesa, mas poucos tratam das normas brasileiras.
% TODO revisar esta seção antes de enviar

\section{Metodologia}\label{sec:metodologia}

A ferramenta segue o fluxo mostrado na Figura~\ref{fig:inexistente}. Primeiro, os
arquivos são descobertos; depois, o texto é extraído e analisado por verificadores
independentes, conforme sugerido por \citeonline{inventado2021}.

\begin{figure}[htb]
  \centering
  \includegraphics[width=0.5\textwidth]{figuras/fluxo.png}
  \label{fig:fluxo}
\end{figure}

A legibilidade é estimada pelo índice de \citeonline{flesch1948} apud
\citeonline{martins1996}, adaptado ao português. A Tabela~\ref{tab:metricas}
resume as métricas utilizadas.

\begin{table}[htb]
  \caption{Métricas de escrita utilizadas}\label{tab:metricas}
  \begin{tabular}{ll}
    Métrica & Faixa \\
    Flesch & 50--75 \\
  \end{tabular}
\end{table}

A pontuação final é dada por
\begin{equation}
  F = 248{,}835 - 1{,}015\,\mathrm{ASL} - 84{,}6\,\mathrm{ASW}.
\end{equation}

\section{Resultados}\label{sec:resultados}

Nos testes realizados com o artigo de exemplo, a ferramenta encontrou 14 dos 14
problemas plantados, em 1,8 segundo, numa máquina comum. Medimos também o tempo
de execução em uma pasta com 20 arquivos, que ficou abaixo de cinco segundos.

É importante ressaltar que a revisão automatizada desempenha um papel crucial no
cenário atual da produção científica. Além disso, vale destacar que a
ferramenta oferece uma abordagem abrangente e robusta para a análise de textos.
Nesse sentido, é fundamental compreender que a qualidade da escrita impacta
diretamente a aceitação dos manuscritos. Dessa forma, a ferramenta se mostra
uma solução inovadora e eficaz para os desafios enfrentados pelos autores.
Em suma, a revisão automatizada representa um avanço significativo no campo.

Os resultados foram analisados pelos autores e foram considerados muito
satisfatórios, e a ferramenta foi avaliada por revisores que consideraram que os
achados eram relevantes, que as sugestões eram claras e que a localização das
linhas estava correta na grande maioria dos casos que foram examinados por eles.

\section{Conclusão}\label{sec:conclusao}

O objetivo deste trabalho é descrever uma ferramenta de auditoria que lê os
arquivos-fonte de um artigo, identifica problemas de estrutura, de referências
e de escrita, e apresenta cada achado com o arquivo, a linha e uma sugestão de
correção, sem jamais alterar os arquivos originais do autor.

Como trabalho futuro, pretende-se ampliar o corpus de comparação e avaliar a
ferramenta com autores reais.


\postextual
\bibliography{referencias}
\end{document}
